\chapter{Esercizi Bilancio}
\section{Esercizi di bilancio}
\subsection{Esercizio: chiusura di bilancio}
\subsubsection*{Testo e dati iniziali}

\noindent Una società presenta un capitale sociale di
20.000, investito inizialmente in:

\begin{itemize}
    \item un terreno del valore di 15.000;
    \item disponibilità in banca per 5.000.
\end{itemize}

\noindent Durante l'anno vengono effettuate dieci
operazioni. Dobbiamo ricostruirle e preparare:

\begin{itemize}
    \item il \textbf{Conto Economico};
    \item lo \textbf{Stato Patrimoniale}.
\end{itemize}

\noindent L'aliquota di imposta è pari al 30\%.

\subsubsection*{Brevissimo richiamo alla teoria}

\noindent Il \textbf{Conto Economico} ci dice
se l'azienda ha guadagnato oppure perso durante l'anno.

\[
\boxed{
\text{Utile prima delle imposte}
=
\text{Ricavi}-\text{Costi}
}
\]

\[
\boxed{
\text{Utile netto}
=
\text{Utile prima delle imposte}-\text{Imposte}
}
\]

\noindent Lo \textbf{Stato Patrimoniale} ci dice
che cosa possiede l'azienda e come è stato finanziato.

\[
\boxed{
\text{Attività}
=
\text{Passività}+\text{Patrimonio netto}
}
\]

\begin{itemize}
    \item \textbf{Attività}: terreno, denaro e crediti.
    \item \textbf{Passività}: debiti verso fornitori
    e debiti tributari.
    \item \textbf{Patrimonio netto}: capitale dei soci
    e utile rimasto nell'azienda.
\end{itemize}

\medskip

\noindent \textbf{Regola fondamentale:}
un costo e il suo pagamento sono due cose diverse.

\noindent Se compriamo oggi e paghiamo domani,
il costo nasce oggi. Domani, quando paghiamo,
eliminiamo il debito: non registriamo nuovamente il costo.

\noindent Lo stesso vale per i ricavi:
una vendita genera un ricavo anche se il cliente
pagherà successivamente.

\medskip

\noindent Per seguire lo svolgimento proposto,
consideriamo tutti i movimenti di denaro nel conto
\textbf{Banca}, anche quando il testo parla di contanti.
Non sono indicate IVA, rimanenze finali o altre
rettifiche: non le introduciamo.

\subsubsection*{Situazione iniziale}

\noindent Il capitale sociale indica quanto
i soci hanno conferito all'azienda.

\noindent Non è un ricavo: non deriva da una vendita.

\[
\text{Terreno}+\text{Banca}
=
15.000+5.000=20.000
\]

\[
\boxed{
\text{Attività iniziali}
=
\text{Capitale sociale}
=
20.000
}
\]

\noindent All'inizio non ci sono debiti.

\subsubsection*{1. Acquisto di materie prime per 2.000
con pagamento dilazionato}

\noindent \textbf{Pagamento dilazionato} significa:
compriamo subito, ma pagheremo in seguito.

\noindent Registriamo quindi:

\begin{itemize}
    \item un costo per materie prime di 2.000;
    \item un debito verso fornitori di 2.000;
    \item nessuna uscita immediata di denaro.
\end{itemize}

\[
\text{Debiti verso fornitori}=2.000
\]

\[
\boxed{\text{Banca}=5.000}
\]

\noindent Il costo esiste già, anche se non abbiamo
ancora pagato il fornitore.

\subsubsection*{2. Vendita di prodotti finiti per 14.000
con incasso immediato}

\noindent Vendiamo prodotti e riceviamo subito
il denaro.

\begin{itemize}
    \item Nasce un ricavo di vendita di 14.000.
    \item La banca aumenta di 14.000.
\end{itemize}

\[
\text{Banca}=5.000+14.000
\]

\[
\boxed{\text{Banca}=19.000}
\]

\subsubsection*{3. Pagamento del debito dell'operazione 1}

\noindent Paghiamo i 2.000 che dovevamo al fornitore.

\[
\text{Banca}=19.000-2.000
\]

\[
\boxed{\text{Banca}=17.000}
\]

\[
\text{Debiti verso fornitori}=2.000-2.000
\]

\[
\boxed{\text{Debiti verso fornitori}=0}
\]

\noindent \textbf{Attenzione:} non nasce un nuovo costo.
Il costo delle materie prime era già stato registrato
nell'operazione 1.

\subsubsection*{4. Acquisto di materie prime per 6.000:
un terzo subito e due terzi dilazionati}

\noindent Il costo dell'acquisto è tutto pari a 6.000.
Il pagamento viene invece diviso in due parti.

\medskip

\noindent La parte pagata subito è:

\[
6.000\cdot\frac{1}{3}=2.000
\]

\noindent La parte che pagheremo successivamente è:

\[
6.000\cdot\frac{2}{3}=4.000
\]

\noindent Quindi:

\begin{itemize}
    \item registriamo un costo per materie prime di 6.000;
    \item paghiamo subito 2.000;
    \item registriamo un debito di 4.000.
\end{itemize}

\[
\text{Banca}=17.000-2.000
\]

\[
\boxed{\text{Banca}=15.000}
\]

\[
\boxed{\text{Debiti verso fornitori}=4.000}
\]

\noindent Gli acquisti complessivi di materie prime
sono ora:

\[
2.000+6.000=\boxed{8.000}
\]

\subsubsection*{5. Pagamento di stipendi per 5.000}

\noindent Gli stipendi rappresentano il costo
del lavoro svolto dai dipendenti.

\begin{itemize}
    \item Registriamo un costo per stipendi di 5.000.
    \item La banca diminuisce di 5.000.
\end{itemize}

\[
\text{Banca}=15.000-5.000
\]

\[
\boxed{\text{Banca}=10.000}
\]

\subsubsection*{6. Pagamento della metà del debito
dell'operazione 4}

\noindent Il debito nato nell'operazione 4
era pari a 4.000.

\noindent Ne paghiamo la metà:

\[
4.000\cdot\frac{1}{2}=2.000
\]

\noindent La banca diventa:

\[
\text{Banca}=10.000-2.000
\]

\[
\boxed{\text{Banca}=8.000}
\]

\noindent Il debito ancora da pagare è:

\[
\text{Debiti verso fornitori}=4.000-2.000
\]

\[
\boxed{\text{Debiti verso fornitori}=2.000}
\]

\noindent Anche in questo caso non nasce un nuovo costo:
stiamo soltanto pagando una parte di un acquisto
già registrato.

\subsubsection*{7. Incasso di un canone di locazione
del terreno per 3.000}

\noindent L'azienda concede il terreno in uso
a un'altra persona e riceve un affitto.

\begin{itemize}
    \item Registriamo un ricavo per
    \textbf{fitti attivi} di 3.000.
    \item La banca aumenta di 3.000.
\end{itemize}

\[
\text{Banca}=8.000+3.000
\]

\[
\boxed{\text{Banca}=11.000}
\]

\noindent Il terreno rimane di proprietà dell'azienda:
non è stato venduto.

\[
\boxed{\text{Terreno}=15.000}
\]

\subsubsection*{8. Vendita di prodotti finiti per 6.000
con incasso dilazionato}

\noindent Vendiamo i prodotti, ma il cliente
pagherà successivamente.

\begin{itemize}
    \item Registriamo subito un ricavo di 6.000.
    \item Nasce un credito verso clienti di 6.000.
    \item La banca non cambia.
\end{itemize}

\[
\boxed{\text{Crediti verso clienti}=6.000}
\]

\[
\boxed{\text{Banca}=11.000}
\]

\noindent Un \textbf{credito} è una somma
che abbiamo il diritto di ricevere.

\noindent I ricavi complessivi delle vendite sono:

\[
14.000+6.000=\boxed{20.000}
\]

\subsubsection*{9. Incasso di un terzo del credito
dell'operazione 8}

\noindent Il cliente paga un terzo dei 6.000
che ci doveva:

\[
6.000\cdot\frac{1}{3}=2.000
\]

\noindent La banca aumenta:

\[
\text{Banca}=11.000+2.000
\]

\[
\boxed{\text{Banca}=13.000}
\]

\noindent Il credito diminuisce:

\[
\text{Crediti verso clienti}=6.000-2.000
\]

\[
\boxed{\text{Crediti verso clienti}=4.000}
\]

\noindent \textbf{Attenzione:} non nasce un altro ricavo.
Il ricavo era già stato registrato al momento
della vendita, nell'operazione 8.

\subsubsection*{10. Pagamento di internet e utenze
per 4.000}

\noindent Paghiamo complessivamente 4.000
per internet, luce e gas.

\begin{itemize}
    \item Registriamo un costo per servizi
    e utenze di 4.000.
    \item La banca diminuisce di 4.000.
\end{itemize}

\[
\text{Banca}=13.000-4.000
\]

\[
\boxed{\text{Banca finale}=9.000}
\]

\subsubsection*{Controllo del saldo della banca}

\noindent Per verificare il risultato,
sommiamo separatamente le entrate e le uscite.

\noindent Le entrate dell'anno sono:

\[
\begin{aligned}
\text{Entrate}
&=14.000+3.000+2.000\\
&=\boxed{19.000}
\end{aligned}
\]

\noindent Le uscite dell'anno sono:

\[
\begin{aligned}
\text{Uscite}
&=2.000+2.000+5.000+2.000+4.000\\
&=\boxed{15.000}
\end{aligned}
\]

\noindent Al denaro iniziale aggiungiamo le entrate
e sottraiamo le uscite:

\[
\begin{aligned}
\text{Banca finale}
&=5.000+19.000-15.000\\
&=\boxed{9.000}
\end{aligned}
\]

\subsubsection*{Preparazione del Conto Economico}

\noindent Nel Conto Economico inseriamo
i ricavi e i costi dell'anno.

\noindent Non inseriamo il capitale sociale,
il terreno, i crediti o i debiti:
questi appartengono allo Stato Patrimoniale.

\medskip

\noindent \textbf{1. Sommiamo i ricavi}

\[
\text{Ricavi di vendita}=14.000+6.000=20.000
\]

\[
\text{Fitti attivi}=3.000
\]

\[
\text{Ricavi totali}=20.000+3.000
\]

\[
\boxed{\text{Ricavi totali}=23.000}
\]

\medskip

\noindent \textbf{2. Sommiamo i costi prima delle imposte}

\[
\text{Materie prime}=2.000+6.000=8.000
\]

\[
\text{Stipendi}=5.000
\]

\[
\text{Internet e utenze}=4.000
\]

\[
\text{Costi totali}=8.000+5.000+4.000
\]

\[
\boxed{\text{Costi prima delle imposte}=17.000}
\]

\medskip

\noindent \textbf{3. Calcoliamo l'utile prima delle imposte}

\[
\begin{aligned}
\text{Utile prima delle imposte}
&=\text{Ricavi}-\text{Costi}\\
&=23.000-17.000\\
&=\boxed{6.000}
\end{aligned}
\]

\medskip

\noindent \textbf{4. Calcoliamo le imposte}

\noindent In questo esercizio applichiamo
il 30\% all'utile prima delle imposte.

\[
\boxed{
\text{Imposte}
=
\text{Utile prima delle imposte}
\cdot\text{Aliquota}
}
\]

\[
30\%=\frac{30}{100}=0{,}30
\]

\[
\text{Imposte}=6.000\cdot0{,}30
\]

\[
\boxed{\text{Imposte}=1.800}
\]

\noindent Le imposte sono un costo dell'anno.
Poiché il testo non ne indica il pagamento,
registriamo anche un debito tributario di 1.800.

\noindent La banca resta quindi pari a 9.000.

\medskip

\noindent \textbf{5. Calcoliamo l'utile netto}

\[
\begin{aligned}
\text{Utile netto}
&=\text{Utile prima delle imposte}-\text{Imposte}\\
&=6.000-1.800\\
&=\boxed{4.200}
\end{aligned}
\]

\subsubsection*{Conto Economico finale}

\begin{center}
    \begin{tabular}{|l|r|}
        \hline
        \textbf{Voce} & \textbf{Importo} \\
        \hline
        Ricavi di vendita & 20.000 \\
        Fitti attivi & 3.000 \\
        \hline
        \textbf{Totale ricavi} & \textbf{23.000} \\
        \hline
        Acquisti di materie prime & 8.000 \\
        Stipendi & 5.000 \\
        Internet e utenze & 4.000 \\
        \hline
        \textbf{Totale costi prima delle imposte}
        & \textbf{17.000} \\
        \hline
        \textbf{Utile prima delle imposte}
        & \textbf{6.000} \\
        Imposte dell'esercizio & 1.800 \\
        \hline
        \textbf{Utile netto} & \textbf{4.200} \\
        \hline
    \end{tabular}
\end{center}

\subsubsection*{Preparazione dello Stato Patrimoniale}

\noindent Adesso fotografiamo la situazione
dell'azienda alla fine dell'anno.

\medskip

\noindent \textbf{1. Attività: che cosa possiede l'azienda?}

\begin{itemize}
    \item Terreno: 15.000.
    \item Crediti ancora da incassare: 4.000.
    \item Denaro in banca: 9.000.
\end{itemize}

\[
\text{Totale attività}=15.000+4.000+9.000
\]

\[
\boxed{\text{Totale attività}=28.000}
\]

\medskip

\noindent \textbf{2. Passività: quanto deve ancora pagare?}

\begin{itemize}
    \item Debiti verso fornitori: 2.000.
    \item Debiti tributari: 1.800.
\end{itemize}

\[
\text{Totale passività}=2.000+1.800
\]

\[
\boxed{\text{Totale passività}=3.800}
\]

\medskip

\noindent \textbf{3. Patrimonio netto: qual è la parte
che appartiene ai soci?}

\noindent Al capitale sociale iniziale
aggiungiamo l'utile netto dell'anno,
che non risulta distribuito ai soci.

\[
\boxed{
\text{Patrimonio netto}
=
\text{Capitale sociale}+\text{Utile netto}
}
\]

\[
\text{Patrimonio netto}=20.000+4.200
\]

\[
\boxed{\text{Patrimonio netto}=24.200}
\]

\subsubsection*{Stato Patrimoniale finale}

\begin{center}
    \begin{tabular}{|l|r|}
        \hline
        \textbf{Attività} & \textbf{Importo} \\
        \hline
        \textbf{Immobilizzazioni materiali} & \\
        Terreno & 15.000 \\
        \hline
        \textbf{Attivo circolante} & \\
        Crediti verso clienti & 4.000 \\
        Banca & 9.000 \\
        \hline
        \textbf{Totale attività} & \textbf{28.000} \\
        \hline
    \end{tabular}
\end{center}

\begin{center}
    \begin{tabular}{|l|r|}
        \hline
        \textbf{Patrimonio netto e passività}
        & \textbf{Importo} \\
        \hline
        \textbf{Patrimonio netto} & \\
        Capitale sociale & 20.000 \\
        Utile dell'esercizio & 4.200 \\
        \textbf{Totale patrimonio netto} & \textbf{24.200} \\
        \hline
        \textbf{Passività} & \\
        Debiti verso fornitori & 2.000 \\
        Debiti tributari & 1.800 \\
        \textbf{Totale passività} & \textbf{3.800} \\
        \hline
        \textbf{Totale patrimonio netto e passività}
        & \textbf{28.000} \\
        \hline
    \end{tabular}
\end{center}

\subsubsection*{Verifica finale}

\noindent Lo Stato Patrimoniale deve rispettare
questa uguaglianza:

\[
\text{Attività}
=
\text{Patrimonio netto}+\text{Passività}
\]

\[
28.000=24.200+3.800
\]

\[
\boxed{28.000=28.000}
\]

\noindent I due totali coincidono:
lo Stato Patrimoniale quadra.

\medskip

\noindent \textbf{Da ricordare:}
l'utile netto è 4.200, mentre il denaro finale
in banca è 9.000.

\noindent Sono valori diversi perché
\textbf{guadagnare e incassare non sono la stessa cosa}:
alcune vendite devono ancora essere incassate
e alcuni costi devono ancora essere pagati.